\chapter{Relativsätze}
\label{cha:Relativ}

Relativsätze werden durch Relativpronomen eingeleitet und beschreiben ein Nomen näher.

\section{Relativpronomen}
\label{sec:Relativpronomen}

Die Relativpronomen richten sich nach Genus, Numerus und Kasus des Bezugsworts:

\begin{center}
\begin{tabular}{|l|c|c|c|c|}
\hline
\textbf{Kasus} & \textbf{Maskulin} & \textbf{Feminin} & \textbf{Neutral} & \textbf{Plural} \\ \hline
Nominativ & der & die & das & die \\ \hline
Akkusativ & den & die & das & die \\ \hline
Dativ & dem & der & dem & denen \\ \hline
Genitiv & dessen & deren & dessen & deren \\ \hline
\end{tabular}
\end{center}

\section{Relativpronomen mit Präpositionen}
\label{sec:RelativPraep}

\begin{center}
\begin{tabular}{|l|l|}
\hline
\textbf{Präposition} & \textbf{Beispiel} \\ \hline
an + dem & an dem = an dem \\ \hline
in + dem & in dem = in dem \\ \hline
auf + das & auf das = auf das \\ \hline
mit + dem & mit dem = mit dem \\ \hline
über + das & über das = über das \\ \hline
\end{tabular}
\end{center}

\section{Übungen zu Relativsätzen}
\label{sec:RelativUebungen}

\subsection{Relativpronomen (Nom/Akk/Dat)}
\begin{enumerate}
    \item Das ist das Buch, \underline{\hspace{1cm}} mich interessiert. \quad \textit{(Lösung: das)}
    \item Der Mann, \underline{\hspace{1cm}} dort steht, ist mein Vater. \quad \textit{(Lösung: der)}
    \item Die Frau, \underline{\hspace{1cm}} singt, ist berühmt. \quad \textit{(Lösung: die)}
    \item Das ist der Mann, \underline{\hspace{1cm}} ich kenne. \quad \textit{(Lösung: den)}
    \item Die Frau, \underline{\hspace{1cm}} ich sah, war berühmt. \quad \textit{(Lösung: die)}
    \item Der Lehrer, \underline{\hspace{1cm}} ich Fragen stellte, antwortete. \quad \textit{(Lösung: dem)}
    \item Die Nachbarn, \underline{\hspace{1cm}} ich grüße, sind freundlich. \quad \textit{(Lösung: denen)}
    \item Das Land, \underline{\hspace{1cm}} ich entstamme, ist Griechenland. \quad \textit{(Lösung: dem)}
\end{enumerate}

\subsection{Relativsätze mit Präpositionen}
\begin{enumerate}
    \item Ich habe ein Problem. Ich spreche mit dem Lehrer. $\rightarrow$ Relativsatz \quad \textit{(Lösung: Ich habe ein Problem, mit dem ich mit dem Lehrer spreche.)}
    \item Sie arbeitet in einer Firma. Sie ist stolz darauf. $\rightarrow$ Relativsatz \quad \textit{(Lösung: Sie arbeitet in einer Firma, auf die sie stolz ist.)}
    \item Er hat einen Freund. Er vertraut ihm. $\rightarrow$ Relativsatz \quad \textit{(Lösung: Er hat einen Freund, dem er vertraut.)}
\end{enumerate}

\chapterend